\section{The recovered source and the original pointwise theorem}
\label{sec:v104-scalar}
Retain the original schedule in \eqref{eq:v104-schedule}. Let
$H_B$ be exactly the original selector for the source which fails
all three preceding clocks, so that
$\mathsf U_B^{\rm rem}=T_{H_B}^{\varepsilon(B)}$. This selector
includes the original short-source test and is supported on
$\ell\le\delta(B)$. It is not redefined by the present comparison.

Choose comparison pairs as in Theorem~\ref{thm:v104-block}, with
$r=\delta$, $\Delta=4\delta$ and constants $N_B,\kappa_B$ fixed
before the collision limit. They may, for example, be the disjoint
critical arms of Theorem~\ref{thm:v104-stationary}. The exact split is
\begin{align}
 \mathsf V_B^{\rm int}(t)
     &=\sum_j\alpha_jT_{H_B,j}^{\varepsilon(B)}(t),\notag\\
 \mathsf R_B^{\rm int}(t)
     &=T_{H_B\one_{(\cup_jI_j)^c}}^{\varepsilon(B)}(t)
                  +\sum_j(1-\alpha_j)T_{H_B,j}^{\varepsilon(B)}(t).
                                                   \label{eq:v104-source-split}
\end{align}
Hence $\mathsf U_B^{\rm rem}=\mathsf V_B^{\rm int}+
\mathsf R_B^{\rm int}$ at finite count, with both terms nonnegative.
The second sum is an explicitly weighted shortage of comparison
trace, with upper bound \eqref{eq:v104-deficit}. The uncovered source
in the first term is retained in full.

\begin{theorem}[Positive recovery at the original outer schedule]
\label{thm:v104-recovery}
Suppose $N_B\kappa_B\le C\varepsilon(B)^{-1/32}$ uniformly in
collision count, radius and exact labels. Then
\begin{equation}\label{eq:v104-paid-recovery}
 \mathcal H_M(\mathsf V_B^{\rm int})\le C_MB^{-1/384},
 \qquad
 \mathcal E_M\le\mathcal H_M(\mathsf R_B^{\rm int})
                                      +C_MB^{-1/384}.
\end{equation}
If $N_B\kappa_B$ is bounded independently of $B$, the first bound
improves to $C_MB^{-1/192}$. Along the same schedule the original
pointwise assertion is equivalent to
$\mathcal H_M(\mathsf R_B^{\rm int})\to0$.
\end{theorem}
\begin{proof}
Apply Corollary~\ref{cor:v104-block-height} to the exact starting
selector $H_B$. Since $\delta=O(\varepsilon^4)$,
\[
 \mathcal H_M(\mathsf V_B^{\rm int})
 \le C N_B\kappa_B
       \{\varepsilon^{1/16}+(1+\varepsilon^{-1})4\delta\}
 \le C(\varepsilon^{1/32}+\varepsilon^{95/32}).
\]
This is $O(B^{-1/384})$; for bounded $N_B\kappa_B$ the leading
power is $\varepsilon^{1/16}=O(B^{-1/192})$. Insert the exact
positive split \eqref{eq:v104-source-split} into the inherited scalar
bound \eqref{eq:v104-old-remainder}. This proves sufficiency.
Conversely, the retained necessity statement gives
$\mathcal H_M(\mathsf U_B^{\rm rem})\to0$ under scalar closure.
Since $0\le\mathsf R_B^{\rm int}\le\mathsf U_B^{\rm rem}$,
necessity follows. The complete physical sums precede the height
functional in every step.
\end{proof}

For the critical-arm theorem a uniform $K_0$ and $\mu_0>0$ give
the stronger rate for the entire selected family. More generally,
any outer choice with $K_B/\mu_B\le C\varepsilon(B)^{-1/32}$
gives the first rate. This allows an expanding family of finite-order
profiles, but is not a claim that their order or trace fraction is
uniformly controlled on all physical pairs. Every such threshold is
chosen from $B$, never from the collision count.

\begin{corollary}[Removing the whole stationary-critical arm contribution]
\label{cor:v104-whole-arm-recovery}
Suppose the comparison pairs are the disjoint actual critical arms
of Theorem~\ref{thm:v104-paid-arms}, with
$\kappa_B\le C\varepsilon(B)^{-1/32}$. Then
\[
 \mathcal H_M\!\left(\sum_j(1-\alpha_j)
                 T_{H_B,j}^{\varepsilon(B)}\right)
                  \le C\varepsilon(B)^{95/32}.
\]
The complete original source on these arms, including its entire
three-clock-failure part, has height $O(B^{-1/384})$. With a bounded
geometric ratio the latter rate is $O(B^{-1/192})$. In particular,
setting $\mathsf U_B^{\rm outside}
=T_{H_B\one_{(\cup_jI_j)^c}}^{\varepsilon(B)}$, one has
\begin{equation}\label{eq:v104-outside-scalar}
 \mathcal E_M\le\mathcal H_M(\mathsf U_B^{\rm outside})
                                      +C_MB^{-1/384}.
\end{equation}
Along this schedule the original pointwise assertion is equivalent
to $\mathcal H_M(\mathsf U_B^{\rm outside})\to0$.
\end{corollary}
\begin{proof}
Apply \eqref{eq:v104-paid-deficit} on the original starting selector.
The complete uncut volume estimate gives
$C\kappa_B(1+\varepsilon^{-1})\delta
=O(\varepsilon^{95/32})$. Adding this paid shortage to
\eqref{eq:v104-paid-recovery} removes the entire angular-arm source,
not only a fraction admitted by trace balance. The exact decomposition
into that source and its positive outside complement proves
\eqref{eq:v104-outside-scalar}. Necessity follows from domination by
the original $\mathsf U_B^{\rm rem}$, and sufficiency from the
new estimate. No count-dependent width is chosen.
\end{proof}

\subsection{The remaining physical estimate}
The new comparison has an actual essential-height conclusion on the
positive recovered source. It does not turn an average, a signed
boundary current or a finite local norm into the target theorem.
Stationary endpoint speed and rapid pointwise variation of the minimum
are no longer obstructions on the interval families just estimated.
The model calculation shows that this recovered family need not have
been admitted by any earlier clock.

Two separate matters are still not proved on the complete Lorentz
family: coverage by comparison intervals with a schedule-compatible
physical overlap and mass balance, and vanishing of the uncovered source. On the
critical arms of Theorem~\ref{thm:v104-paid-arms}, the entire
comparison-trace shortage is already paid by the original uncut volume. In particular a flat endpoint
pulse on an entire isolated chamber need not have any comparison
interval satisfying \eqref{eq:v104-inclusion}; its positive mass is
not removed. Loss of physical continuation before the comparison
levels, and singular critical extensions remain in the
outside source. General comparison families without the critical
incoming normalization may still have a trace shortage in
$\mathsf R_B^{\rm int}$.

Thus \eqref{eq:v104-paid-recovery} is not a proof that the full
remaining source has vanishing height. The exact-label pointwise
theorem remains the governing assertion \eqref{eq:v76-target}; its
full closure and the corresponding unrestricted point-conditioned
consequences are not promoted. The original arithmetic zero classes
remain in the unnormalized theorem, and scalar closure still precedes
any denominator, same-roof ratio or unrestricted likelihood statement.
